Laser ultrasonic testing offers significant advantages over conventional contact transducers: it enables non-contact inspection, supports high-speed scanning, and can be applied to complex geometries and rough surfaces without coupling agents \cite{liu2022deep}. These properties make laser ultrasound particularly attractive for non-destructive testing in aerospace, manufacturing, and process control environments. However, laser-generated ultrasonic signals suffer from low signal-to-noise ratio (SNR) due to the stochastic nature of the laser ablation process and thermoelastic conversion \cite{zhang2020deep}. The resulting noise obscures diagnostically relevant acoustic features---including wave arrival times, amplitude ratios, and waveform morphology---that are essential for reliable defect detection and characterization \cite{chen2019wavelet}.

Signal averaging is the most widely adopted method for SNR improvement in laser ultrasound \cite{zhang2020deep}. By acquiring and summing multiple repetitions of the same measurement, random noise components cancel out while the coherent signal builds up, yielding an SNR gain proportional to the square root of the number of averages. Despite its effectiveness, extensive averaging imposes practical constraints: it increases acquisition time linearly, can induce surface damage on sensitive or fragile materials under repeated laser pulses, and becomes impractical for high-speed industrial inspection \cite{yang2023unsupervised}. These limitations motivate the need for denoising methods that can achieve comparable or superior SNR improvement without relying on large numbers of averages.

Conventional denoising techniques---including Wiener filtering, wavelet thresholding, and block-matching collaborative filtering (BM3D)---can suppress noise to some extent \cite{dabov2007bm3d, chen2019wavelet}. However, these methods typically optimize for aggregate noise metrics such as mean squared error or SNR improvement, without explicit constraints on the preservation of acoustically meaningful features \cite{yang2023unsupervised}. For defect detection, the fidelity of diagnostic features---the exact arrival time of dispersive modes, relative amplitudes of reflected echoes, and the shape of the waveform envelope---is as important as overall noise suppression. Suppressing noise while degrading these features provides little practical benefit. This work addresses the problem of achieving effective laser ultrasonic signal denoising that explicitly preserves acoustic features critical for defect detection, thereby reducing dependence on extensive signal averaging.
